% -------------------------------------------------------------
% PHẦN VI: KẾT LUẬN \& ĐỀ XUẤT PHÁT TRIỂN
% -------------------------------------------------------------
\section{Kết luận \& Hướng phát triển}

\begin{frame}{27. Kết quả Đạt được \& Đóng góp Khoa học}
    \textbf{Đóng góp Kỹ thuật:}
    \begin{enumerate}
        \item Khung mô hình hoá dự án \textbf{Heterogeneous Graph} với 3 Node Types (Task/Resource/Shift).
        \item Áp dụng thành công \textbf{Masked Autoencoder SSL} (Phase 0) giải quyết bài toán thiếu nhãn.
        \item Khai sinh khái niệm \textbf{AI-augmented Beta-PERT}: AI hiệu chỉnh tham số $(a, m, b)$ riêng biệt.
        \item Lập lịch 4-Mode với \textbf{CP-SAT AI-Guided Pareto} giải quyết bài toán đa mục tiêu.
        \item Phát triển hệ thống phần mềm Full-Stack (React/FastAPI) \textbf{có tính ứng dụng thực tiễn cao}.
    \end{enumerate}
    
    \vspace{0.2cm}
    \textbf{Đóng góp Học thuật:} 
    Kết hợp 4 kỹ thuật (HGTConv + Masked Autoencoder + MC-PERT + CP-SAT) trong 1 Pipeline liền mạch.
\end{frame}

\begin{frame}{28. Đánh giá Ưu \& Nhược điểm của Mô hình AI}
    \begin{columns}
        \column{0.5\textwidth}
        \textbf{Ưu điểm kỹ thuật:}
        \begin{itemize}
            \item \textbf{Cải thiện độ chính xác (Mean R² = 0.3801):} Đạt mức hiệu suất cao hơn đáng kể so với mô hình Baseline tĩnh.
            \item \textbf{Khả năng xử lý đồ thị quy mô lớn:} Đạt hiệu quả cao trên các mạng lưới phức tạp (VD: đồ thị >400 đỉnh, R² = 0.82).
            \item \textbf{Khung huấn luyện linh hoạt:} Tích hợp quy trình tự động (Pretrain $\rightarrow$ Transductive $\rightarrow$ Optuna HPO) giúp rút ngắn thời gian tinh chỉnh tham số.
        \end{itemize}
        
        \column{0.5\textwidth}
        \textbf{Hạn chế của hệ thống:}
        \begin{itemize}
            \item \textbf{Hạn chế về lượng mẫu dữ liệu:} Kích thước tập mẫu nhỏ ($N=5$) buộc phải áp dụng học chuyển nạp (Transductive Learning), giới hạn khả năng học quy nạp (Inductive Learning) trên các đồ thị mới.
            \item \textbf{Hiệu năng giảm trên đồ thị thưa (Sparse Graph):} Các đồ thị $<50$ đỉnh hạn chế quá trình truyền thông điệp (Message Passing), ảnh hưởng đến mức độ hội tụ.
            \item \textbf{Đồ thị tĩnh (Static Graph):} Dữ liệu hiện tại giới hạn ở mức biểu diễn không gian tĩnh, chưa mô hình hóa được diễn biến theo chuỗi thời gian thực.
        \end{itemize}
    \end{columns}
\end{frame}

\begin{frame}{29. Định hướng nghiên cứu tiếp theo}
    Từ các hạn chế hiện tại, một số hướng tiếp cận được đề xuất:
    \vspace{0.2cm}
    \begin{table}
        \centering
        \footnotesize
        \begin{tabular}{@{}llp{6.5cm}@{}}
            \toprule
            \textbf{Vấn đề} & \textbf{Định hướng phát triển} & \textbf{Giải pháp Kỹ thuật} \\
            \midrule
            \textbf{Cấu trúc đồ thị thưa} & Tối ưu cơ chế truyền tin & Bổ sung cơ chế \textbf{Virtual Nodes} nhằm giảm độ dài đường đi trung bình và tăng hiệu suất Message Passing. \\
            \midrule
            \textbf{Hạn chế tập dữ liệu} & Tăng cường tính khái quát hóa & Mở rộng quy mô tập dữ liệu (>100 dự án) nhằm hỗ trợ phương pháp học quy nạp (Inductive Learning) cho bài toán Zero-shot. \\
            \midrule
            \textbf{Đồ thị không gian tĩnh} & Biểu diễn chuỗi thời gian & Tích hợp \textbf{Temporal HGT (T-HGT)} để mô hình hóa các biến động tiến độ và tài nguyên theo trục thời gian thực. \\
            \bottomrule
        \end{tabular}
    \end{table}
\end{frame}
